% Traced rails and open string tips use the picture frame's boundary.
\documentclass[border=5pt]{standalone}
\usepackage{tenkz}
\makeatletter
\ExplSyntaxOn
\file_input:n { tenkz-kernel.code.tex }
\ExplSyntaxOff
\makeatother
\begin{document}
\tenkzkernel
\begin{tenkz}[
    rows={wire,wire}, cols=3, bonds=none,
    west=trace, east=trace]
  \tn[at=(1,1), skin=box]{A}
  \tn[at=(1,3), skin=box]{B}
  \tn[at=(2,1), skin=box]{C}
  \tn[at=(2,3), skin=box]{D}
\end{tenkz}
\qquad
\begin{tenkz}[
    rows={wire,wire}, cols=3, bonds=none,
    north=trace, south=trace]
  \tn[at=(1,1), skin=box]{A}
  \tn[at=(1,2), skin=box]{B}
  \tn[at=(1,3), skin=box]{C}
  \tn[at=(2,1), skin=box]{D}
  \tn[at=(2,2), skin=box]{E}
  \tn[at=(2,3), skin=box]{F}
\end{tenkz}
\qquad
\begin{tenkz}[rows={wire,wire,wire}, cols=3, bonds=none]
  \tn[at=(1,1), skin=box]{A}
  \tn[at=(1,2), skin=box]{B}
  \tn[at=(1,3), skin=box]{C}
  \tn[at=(2,1), skin=none, name=jH1]{}
  \tn[at=(2,3), skin=none, name=jH2]{}
  \tnwire[kind=string, name=segH1]{open w}{jH1}
  \tnwire[kind=string, name=segH2]{jH1}{jH2}
  \tnwire[kind=string, name=segH3]{jH2}{open e}
\end{tenkz}
\qquad
\begin{tenkz}[rows={wire,wire,wire}, cols=3, bonds=none]
  \tn[at=(1,1), skin=box]{A}
  \tn[at=(2,1), skin=box]{B}
  \tn[at=(3,1), skin=box]{C}
  \tn[at=(1,2), skin=none, name=jV1]{}
  \tn[at=(3,2), skin=none, name=jV2]{}
  \tnwire[kind=string, name=segV1]{open n}{jV1}
  \tnwire[kind=string, name=segV2]{jV1}{jV2}
  \tnwire[kind=string, name=segV3]{jV2}{open s}
\end{tenkz}
\qquad
\end{document}
